%% ma: L12-B08-0006
%% lop: 12
%% chuong: 2
%% bai: 8
%% dang: 12.8.7
%% loai: TN4
%% muc: NB
%% dap_an: "B"
%% nguon: "(NBV)-ĐỀ SỐ 3-MỖI NGÀY 1 ĐỀ THI 2025.pdf (D03, MỖI NGÀY 1 ĐỀ - Đề số 3), Phần I Câu 10"
%% kien_thuc: "Tích có hướng của hai vectơ trong không gian (ngoài SGK KNTT)"
%% trang_thai: cach-ly
%% ghi_chu: "Đáp án gốc B đúng: [a,b]=(-1;4;-3) (kiểm bằng sympy). Dạng 12.8.7 là dạng đề xuất cho câu vượt chương trình"
%% ly_do: "Dùng tích có hướng của hai vectơ, kiến thức không có trong SGK Toán 12 Kết nối tri thức. Gợi ý sửa: đổi thành hỏi tích vô hướng của hai vectơ (Bài 8), hoặc nhờ tìm một vectơ vuông góc với cả hai vectơ đã cho bằng hệ phương trình"
\begin{ex}
Trong không gian $Oxyz$, cho các vectơ $\vec a=(-1;2;3)$, $\vec b=(2;-1;-2)$. Khi đó $[\vec a,\vec b]$ có toạ độ là
\choice{$(1;-4;3)$}{\True $(-1;4;-3)$}{$(-7;8;5)$}{$(7;-8;-5)$}
\loigiai{$[\vec a,\vec b]=\left(2\cdot(-2)-3\cdot(-1);\,3\cdot2-(-1)\cdot(-2);\,(-1)\cdot(-1)-2\cdot2\right)=(-1;4;-3)$.}
\end{ex}
